\section{Filtering, sorting, and duplicate removal}\label{sec:filtering}

Filtering turns a loaded set of records into a smaller, ordered set. Most steps
only read the records; renaming component codes rewrites a record's
composition. No step changes the input database on disk. A filtering run writes
the records that remain with \flg{--out}, prints them with print flags, or both.
Print flags alone do not filter: a print-only run shows every loaded record in
file order.

\subsection{Order of the filtering steps}\label{sec:pipeline-order}

The steps always run in this fixed order, each within the current composition
groups:

\begin{enumerate}
  \item drop records whose status is not \fld{ok};
  \item rename component codes with \flg{--rename}, in the order the flags were
        given;
  \item \flg{--extract} or \flg{--nextract};
  \item cutoffs, in the order the cutoff flags were given;
  \item \flg{--uniq};
  \item \flg{--reacted} or \flg{--reacted2};
  \item \flg{--sort}, then \flg{--reverse};
  \item \flg{--select};
  \item \flg{--arbalign};
  \item \flg{--glob};
  \item \flg{--shuffle}.
\end{enumerate}

The non-\fld{ok} count is reported as a warning:
\texttt{excluded <k> records with status failed or truncated}.

When \flg{--uniq} is used with a list of properties, or when \flg{--select}
or \flg{--arbalign} is used, and no \flg{--sort} is given, the records are
first sorted by \fld{el} ascending. \flg{--uniq dup} does not trigger this
automatic sort.

\subsection{Grouping}\label{sec:grouping}

Most steps work per \emph{composition group}: records with the same
composition form a group, and records without composition information form
one separate unnamed group. Groups are processed in first-appearance order;
within a step, the overall record order is preserved.

\flg{--noname} collapses everything into one global group. Under
\flg{--noname}, \flg{--extract} and \flg{--nextract} compare their patterns
directly against each record's \fld{name} (the same text the \fld{filename}
column prints), with no composition grammar, ranges, wildcards, or zero-count
rules. \flg{--arbalign} is the exception: it always compares records only
when they have the same stored composition, so \flg{--noname} does not make
different compositions comparable.

\subsection{Renaming component codes}\label{sec:rename}

\flg{--rename OLD NEW} moves the count of one component code to another code.
Both codes are plain letters, and upper and lower case are different codes. The
flag may be repeated; renames apply in the order given and chain together, so
\flg{--rename p a --rename a sa} first moves \texttt{p} counts to \texttt{a}
and then moves those \texttt{a} counts on to \texttt{sa}.

The rename applies to every record that has a stored composition:

\begin{itemize}
  \item If the record does not contain \texttt{OLD}, its composition is left
        unchanged.
  \item If \texttt{OLD} and \texttt{NEW} are the same code, nothing changes.
  \item Otherwise the count of \texttt{OLD} is added to the count of
        \texttt{NEW}. For example, a record with 3 \texttt{sa}, 1 \texttt{am},
        and 1 \texttt{p} becomes 4 \texttt{sa} and 1 \texttt{am} under
        \flg{--rename p sa}.
  \item A record with no stored composition is left unchanged.
\end{itemize}

Renaming changes the stored composition, not the record's file name or source
path. The \fld{filename} column still prints the original file stem, while the
\fld{name} and \fld{components} columns and all composition groups show the new
composition. Because renaming is the first filtering step after failed records
are dropped, later steps such as \flg{--extract} and relative cutoffs work with
the new codes.

\subsection{Extracting compositions}\label{sec:extract}

\flg{--extract QUERY} keeps records whose composition matches the query;
\flg{--nextract QUERY} keeps the records that do not. A query is a
comma-separated list of patterns; a record matches if any pattern matches.

A pattern is built from these pieces:

\begin{verbatim}
query   := pattern ( "," pattern )*
pattern := ( pair | range | choice | bare-star | count-star
             | count-star-code | star-code )+
pair    := count code
range   := count "-" count code          (inclusive count expansion)
choice  := "(" count ( "," count )* ")" code
bare-star       := "*"
count-star      := count "*"
count-star-code := count "*" code
star-code       := "*" code
count   := digit+ (decimal, at most 4294967295)
code    := letter+ (plain letters, case matters)
\end{verbatim}

A pair \texttt{(count, code)} with a positive count requires that exact count
for the code. A count of zero requires the code to be absent. Codes the
pattern does not mention are \emph{extras}, governed by the wildcard:

\begin{center}
\begin{tabular}{@{}ll@{}}
\toprule
\textbf{Form} & \textbf{Extra records allowed} \\
\midrule
no wildcard & none: the pattern is an exact composition match \\
\texttt{*} & any composition \\
\texttt{N*} & extras whose counts sum to exactly $N$ \\
\texttt{*code} & only the named codes, at any counts \\
\texttt{N*code} & only the named codes, with counts summing to exactly $N$ \\
\bottomrule
\end{tabular}
\end{center}

Examples:

\begin{center}
\begin{tabular}{@{}ll@{}}
\toprule
\textbf{Query} & \textbf{Matches} \\
\midrule
\texttt{2am3sa} & exactly 2 \texttt{am} and 3 \texttt{sa}, no extras \\
\texttt{2am3sa,1am1sa} & either of the two compositions \\
\texttt{1am0-2w} & \texttt{1am}, \texttt{1am1w}, \texttt{1am2w} \\
\texttt{(2,3)am1w} & \texttt{2am1w} or \texttt{3am1w} \\
\texttt{2am*w} & 2 \texttt{am} plus any number of \texttt{w} \\
\texttt{2am2*w} & 2 \texttt{am} plus \texttt{w} counts summing to 2 \\
\bottomrule
\end{tabular}
\end{center}

Further rules:

\begin{itemize}
  \item A range expands inclusively. A range with more than 1000 count values,
        or a pattern expanding beyond 10000 patterns, is a usage error.
  \item A range whose low end exceeds its high end is a usage error.
  \item Anything after a \texttt{-} that a range did not consume is treated
        as a label comment and ignored, so a full JKQC file stem can be
        pasted directly.
  \item An underscore is not a valid pattern character; a pattern containing
        one is a usage error.
  \item Records without \fld{components} behave as if their composition were
        empty: they match only patterns with no pairs, such as \texttt{*}.
\end{itemize}

\subsection{Cutoffs}\label{sec:cutoffs}

Each cutoff compares a record's property value with a number. The ten flags
are:

\begin{center}
\begin{tabular}{@{}lll@{}}
\toprule
\textbf{Absolute flag} & \textbf{Comparison} & \textbf{Mode} \\
\midrule
\flg{--cut} & $\le$ & absolute \\
\flg{--filter-lt} & $<$ & absolute \\
\flg{--pass} & $>$ & absolute \\
\flg{--filter-ge} & $\ge$ & absolute \\
\flg{--is} & $=$ & absolute \\
\flg{--isnot} & $\ne$ & absolute \\
\midrule
\flg{--cutr} & $\le$ & relative \\
\flg{--rel-filter-lt} & $<$ & relative \\
\flg{--passr} & $>$ & relative \\
\flg{--rel-filter-ge} & $\ge$ & relative \\
\bottomrule
\end{tabular}
\end{center}

They are written \flg{--cut PROPERTY VALUE} and so on, and they may repeat;
cutoffs chain in command order.

\begin{center}
\begin{tabular}{@{}lll@{}}
\toprule
\textbf{Property} & \textbf{Value} & \textbf{Units} \\
\midrule
\fld{el} & electronic energy & hartree absolute; kcal/mol relative \\
\fld{g} & Gibbs free energy & hartree absolute; kcal/mol relative \\
\fld{elcorr} & attached correction energy & hartree absolute; kcal/mol relative \\
\fld{gcorr} & Gibbs from the correction & hartree absolute; kcal/mol relative \\
\fld{rg} & radius of gyration & \AA{} \\
\fld{lf} & lowest stored frequency & cm$^{-1}$ \\
\fld{dip} & dipole magnitude & atomic units \\
\fld{mass} & total mass & u \\
\fld{charge} & net charge & counts \\
\fld{multiplicity} & spin multiplicity & counts \\
\fld{natoms} & geometry atom count & counts \\
\fld{nfreq} & stored frequency count & counts \\
\fld{nimag} & negative frequency count & counts \\
\bottomrule
\end{tabular}
\end{center}

In relative mode the group minimum of the property is subtracted first (over
defined values), then the comparison is applied. Thus
\flg{--cutr el 10} keeps records within 10 kcal/mol of their group's lowest
electronic energy. Under \flg{--noname} the minimum is global.

A record whose property is undefined fails $<$, $\le$, $>$, $\ge$, and $=$,
and passes $\ne$. This mirrors the original tool's treatment of missing
values.

\begin{warning}
Absolute energy cutoffs are in hartree and relative energy cutoffs in
kcal/mol, regardless of whether \flg{--unit} is used for printing. A value
such as \flg{--cut el -38.5} is hartree; \flg{--cutr el 10} is kcal/mol.
\end{warning}

\subsection{Sorting, selecting, and shuffling}\label{sec:sorting}

\flg{--sort PROPERTY} sorts ascending; \flg{--reverse} reverses the
orientation. The possible properties are:

\begin{center}
\begin{tabular}{@{}ll@{}}
\toprule
\textbf{Property} & \textbf{Orders by} \\
\midrule
\fld{el} & electronic energy \\
\fld{g} & Gibbs free energy at the run's conditions \\
\fld{elcorr} & attached correction energy \\
\fld{gcorr} & Gibbs free energy from the correction \\
\fld{rg} & radius of gyration \\
\fld{mass} & total mass \\
\fld{filename} & the record's \fld{name}; nameless records first \\
\bottomrule
\end{tabular}
\end{center}

A record whose numeric sorting value is undefined sorts last ascending and
first descending. \flg{--reverse} alone, with no sort in effect, changes
nothing but still counts as a manipulation.

\flg{--select N} keeps the first $N$ records of each group after sorting.
Groups shorter than $N$ pass whole. \flg{--glob} keeps, per composition
group, the record with the lowest Gibbs free energy at the run's conditions;
ties keep the first in order. Records with undefined Gibbs are dropped.
Under \flg{--noname}, \flg{--glob} keeps the one global minimum.

\flg{--shuffle} puts the final result in a random order. The order is driven
by a fixed starting value, \flg{--seed} (default 0), so the same seed always
gives the same shuffle.

\subsection{Uniqueness}\label{sec:uniqueness}

There are two ways to remove duplicate structures: by name or by recorded
property values, using \flg{--uniq}, and by three-dimensional geometry, using
\flg{--arbalign} (Section~\ref{sec:arbalign}). \flg{--uniq} removes
redundant structures within each composition group. Both modes are stable:
the first occurrence in the current order survives.

\begin{description}
  \item[\flg{--uniq dup}] Drops records that share a \fld{name} with an earlier
        record. Records without a name are never dropped as duplicates.
  \item[\flg{--uniq PROPERTIES}] Drops records whose rounded set of property
        values matches an earlier record in the same group.
\end{description}

The property list is a comma-separated list of \fld{metric[decimals]} items,
for example \flg{--uniq el2.4,rg}. The decimal count is optional; defaults are
\fld{el} 3, \fld{g} 3, \fld{rg} 2, \fld{mass} 2, and \fld{dip} 1. A fractional
decimal count is accepted. Two values agree at $d$ decimals exactly when

\begin{equation}
\lfloor a \cdot 10^d \rfloor = \lfloor b \cdot 10^d \rfloor,
\end{equation}

with floor toward $-\infty$. Undefined metrics are treated as equal to each
other, so two records that both lack the metric collapse.

Because the automatic sort is by \fld{el} ascending, writing
\flg{--uniq el2.4} without \flg{--sort} keeps the lowest-electronic-energy
member of each group of duplicates. Add \flg{--sort g} first if you want the
lowest-Gibbs representative instead.

\subsection{Removing duplicate geometries}\label{sec:arbalign}

Two conformers can have different names and slightly different energies but
still be the same three-dimensional structure. \flg{--arbalign THRESHOLD}
removes such duplicates by comparing geometries. The threshold is a distance
in \AA{}: a record is dropped when its best-fit root-mean-square distance
(RMSD) to an earlier surviving record of the same composition is strictly
less than the threshold. The RMSD is a size in \AA{} that measures how far
apart two structures are after the best possible rotation and atom matching.

Every structure is first turned so that its principal axes (the natural axes
of its mass distribution) line up. This always happens; there is no separate
flag to turn it on. The comparison then tries different atom orderings,
rotations, and mirror-image arrangements, so the order in which atoms appear
in the file and the molecule's orientation do not matter.

The step runs after \flg{--sort} and \flg{--select} and before
\flg{--glob} and \flg{--shuffle}. When no \flg{--sort} is given, \CJKQC{}
first sorts by \fld{el} ascending, so the lowest-electronic-energy member of
a duplicate set survives. Add \flg{--sort g} first if you want the
lowest-Gibbs member instead.

Further rules:

\begin{itemize}
  \item The comparison is strict: a pair exactly at the threshold is kept.
        A threshold of 0 therefore removes nothing.
  \item A negative, non-numeric, or infinite threshold is a usage error.
  \item Records without a usable geometry are kept and reported on the error
        output as
        \texttt{arbalign: kept <k> records without usable geometry}.
  \item A pair whose structures have different numbers or kinds of atoms is
        not compared.
  \item The comparison is done two records at a time: once a record is
        removed, it cannot remove another record. The first surviving record
        in the current order wins.
  \item \flg{--uniq} and \flg{--arbalign} can be combined. \flg{--uniq}
        runs before the sort; then the sort, \flg{--select}, and
        \flg{--arbalign} follow.
  \item \flg{--parallel} also controls how many geometry comparisons run at
        once. The survivors are the same at every setting.
\end{itemize}

\begin{note}
If you use \flg{--arbalign}, please cite the ArbAlign paper:
B.~Temelso, J.~M. Mabey, T.~Kubota, N.~Appiah-Padi, and G.~C. Shields,
``ArbAlign: A Tool for Optimal Alignment of Arbitrarily Ordered Isomers
Using the Kuhn--Munkres Algorithm,'' \emph{J. Chem. Inf. Model.}
\textbf{57} (2017), 1045--1054,
\href{https://doi.org/10.1021/acs.jcim.6b00546}{doi:10.1021/acs.jcim.6b00546}.
\end{note}

For example, keep the lowest-energy member of each group of structures that
are within 0.5 \AA{} RMSD of each other:

\begin{verbatim}
CJKQC db.jsonl --sort el --arbalign 0.5 \
  --filename --el --out unique.jsonl
\end{verbatim}

\subsection{Reacted and exploded structures}\label{sec:reacted}

Within a composition group, most conformers share a fragment fingerprint. A
structure whose optimization reacted or exploded has a different fingerprint.
\CJKQC{} builds each fingerprint from the geometry and uses the group's most
frequent fingerprint as the consensus.

Two atoms are bonded when their distance is below a pair-specific cutoff:

\begin{center}
\begin{tabular}{@{}lc@{}}
\toprule
\textbf{Pair} & \textbf{Cutoff (\AA{})} \\
\midrule
C--N (either order) & 1.75 \\
N--N & 1.50 \\
S--O or O--S & 1.90 \\
O--N or N--O & 1.90 \\
H--H & 0.80 \\
any other pair involving H & never bonded \\
every other heavy--heavy pair & 2.00 \\
\bottomrule
\end{tabular}
\end{center}

Each connected component's fragment formula is its element symbols sorted
alphabetically and concatenated. A record's signature is the sorted list of
its fragment formulas joined with \texttt{+}. For example, a bonded C--N
fragment plus a lone H is \texttt{CN+H}. A record without a geometry has the
signature \texttt{(none)} and participates like any other signature.

\flg{--reacted} keeps the records whose signature equals the group's
consensus; \flg{--reacted2} keeps exactly the deviants. Ties are resolved in a
fixed order: highest count first, then earliest record order.

\subsection{Writing the survivors}\label{sec:filter-out}

A manipulation run writes its survivors with \flg{--out} and reports:

\begin{verbatim}
filtered 17 records into kept.jsonl
\end{verbatim}

Correction and extrapolation add corrected energies to the records before
filtering starts, so cutoffs, sorts, and sections see those values
(Chapter~\ref{sec:corrections}). The overwrite confirmation and \flg{--fout}
rules are the same as in collection.

\subsection{A worked filtering example}\label{sec:filtering-example}

Keep only one conformer per rounded Gibbs/radius signature for each
composition, discard structures more than 5 kcal/mol above the group minimum,
sort by Gibbs energy, and print the lowest 3 per group:

\begin{verbatim}
CJKQC db.jsonl \
  --cutr g 5 \
  --uniq g2,rg1 \
  --sort g \
  --select 3 \
  --filename --name --g --rg --pop
\end{verbatim}

The steps apply in the fixed order: the relative Gibbs cutoff, then
uniqueness, then the sort, then the per-group limit. Populations are computed
over the final printed set, so they renormalize over the three survivors of
each group.
